\documentclass[12pt]{article}

\usepackage[a4paper,margin=1in]{geometry}
\usepackage{amsmath,amssymb,amsthm}
\usepackage{tikz}
\usepackage{enumitem}
\usepackage{xcolor}
\usepackage{array}

\setlength{\parindent}{0pt}
\setlength{\parskip}{0.75em}

\newtheorem{question}{Question}
\newtheorem*{theory}{Theory to Use}
\newtheorem*{answer}{Answer}

\title{Lecture 13 Practice Questions and Step-by-Step Answers}
\author{MA4034 -- Time Series and Stochastic Processes}
\date{}

\begin{document}

\maketitle

\section*{How to Use This File}

Each question first states the theory or formula needed. Then the answer is solved step by step.

\section{General Birth-Death Queue Balance Equations}

\begin{question}
Consider a birth-death queue with state-dependent arrival and service rates:
\[
\lambda_n=3,\qquad n=0,1,2,\ldots
\]
and
\[
\mu_n=5,\qquad n=1,2,3,\ldots
\]
Find \(p_n\) in terms of \(p_0\), and then find \(p_0\).
\end{question}

\begin{theory}
For a birth-death process,
\[
p_n=
\left(
\frac{\lambda_{n-1}\lambda_{n-2}\cdots \lambda_0}
{\mu_n\mu_{n-1}\cdots \mu_1}
\right)p_0,
\qquad n=1,2,\ldots
\]

If \(\lambda_n=\lambda\) and \(\mu_n=\mu\), then
\[
p_n=\rho^n p_0,
\qquad
\rho=\frac{\lambda}{\mu}.
\]

For an infinite queue, the probabilities must sum to 1:
\[
\sum_{n=0}^{\infty}p_n=1.
\]
\end{theory}

\begin{answer}
Here,
\[
\lambda=3,\qquad \mu=5.
\]

Therefore,
\[
\rho=\frac{\lambda}{\mu}=\frac{3}{5}.
\]

So
\[
p_n=\rho^n p_0
=\left(\frac{3}{5}\right)^n p_0.
\]

Now use
\[
\sum_{n=0}^{\infty}p_n=1.
\]

\[
\sum_{n=0}^{\infty}\left(\frac{3}{5}\right)^n p_0=1.
\]

\[
p_0\sum_{n=0}^{\infty}\left(\frac{3}{5}\right)^n=1.
\]

Using the geometric series formula,
\[
\sum_{n=0}^{\infty}r^n=\frac{1}{1-r},\qquad |r|<1,
\]
we get
\[
p_0\frac{1}{1-\frac{3}{5}}=1.
\]

\[
p_0\frac{1}{\frac{2}{5}}=1.
\]

\[
p_0\frac{5}{2}=1.
\]

\[
p_0=\frac{2}{5}.
\]

Therefore,
\[
\boxed{p_n=\left(\frac{3}{5}\right)^n\frac{2}{5}.}
\]
\end{answer}

\section{M/M/1 Infinite Queue}

\begin{question}
A single-server car wash has arrival rate \(\lambda=4\) cars per hour and service rate \(\mu=6\) cars per hour.
Find:
\begin{enumerate}[label=(\alph*)]
    \item \(\rho\),
    \item \(p_0\),
    \item \(p_3\),
    \item \(L_s\),
    \item \(L_q\).
\end{enumerate}
\end{question}

\begin{theory}
For an \(M/M/1\) infinite queue:
\[
\rho=\frac{\lambda}{\mu}.
\]

Steady state exists if
\[
\rho<1.
\]

The state probabilities are
\[
p_n=(1-\rho)\rho^n,\qquad n=0,1,2,\ldots
\]

The expected number in the system is
\[
L_s=\frac{\rho}{1-\rho}.
\]

The expected number in the queue is
\[
L_q=\frac{\rho^2}{1-\rho}.
\]
\end{theory}

\begin{answer}
\textbf{(a) Utilization}
\[
\rho=\frac{\lambda}{\mu}=\frac{4}{6}=\frac{2}{3}.
\]

\[
\boxed{\rho=\frac{2}{3}.}
\]

\textbf{(b) Probability system is empty}
\[
p_0=1-\rho=1-\frac{2}{3}=\frac{1}{3}.
\]

\[
\boxed{p_0=\frac{1}{3}.}
\]

\textbf{(c) Probability of 3 cars in the system}
\[
p_3=(1-\rho)\rho^3.
\]

\[
p_3=\frac{1}{3}\left(\frac{2}{3}\right)^3
=\frac{1}{3}\cdot \frac{8}{27}
=\frac{8}{81}.
\]

\[
\boxed{p_3=\frac{8}{81}\approx 0.0988.}
\]

\textbf{(d) Expected number in the system}
\[
L_s=\frac{\rho}{1-\rho}
=\frac{\frac{2}{3}}{1-\frac{2}{3}}
=\frac{\frac{2}{3}}{\frac{1}{3}}
=2.
\]

\[
\boxed{L_s=2.}
\]

\textbf{(e) Expected number in the queue}
\[
L_q=\frac{\rho^2}{1-\rho}
=\frac{\left(\frac{2}{3}\right)^2}{\frac{1}{3}}
=\frac{\frac{4}{9}}{\frac{1}{3}}
=\frac{4}{3}.
\]

\[
\boxed{L_q=\frac{4}{3}.}
\]
\end{answer}

\section{Using Little's Formula}

\begin{question}
For an \(M/M/1\) queue, suppose
\[
\lambda=5\text{ customers per hour}
\]
and the expected number in the system is
\[
L_s=10.
\]
Find the expected time a customer spends in the system.
\end{question}

\begin{theory}
Little's formula says
\[
L_s=\lambda_{\text{eff}}W_s.
\]

If every arriving customer can join the system, then
\[
\lambda_{\text{eff}}=\lambda.
\]

Therefore,
\[
W_s=\frac{L_s}{\lambda_{\text{eff}}}.
\]
\end{theory}

\begin{answer}
Since the queue is not stated to reject customers,
\[
\lambda_{\text{eff}}=\lambda=5.
\]

Using Little's formula,
\[
L_s=\lambda_{\text{eff}}W_s.
\]

Substitute:
\[
10=5W_s.
\]

\[
W_s=2.
\]

\[
\boxed{W_s=2\text{ hours}.}
\]

Interpretation: On average, a customer spends \(2\) hours in the whole system, including both waiting and service.
\end{answer}

\section{Finite Capacity Single-Server Queue}

\begin{question}
A small barbershop has one barber and only \(3\) waiting chairs.
Including the customer being served, the maximum number in the system is
\[
N=4.
\]
Customers arrive at rate \(\lambda=4\) per hour and the barber serves at rate \(\mu=4\) per hour.
Find:
\begin{enumerate}[label=(\alph*)]
    \item the steady-state probabilities \(p_0,p_1,p_2,p_3,p_4\),
    \item the probability that an arriving customer cannot enter,
    \item the expected number of customers in the system.
\end{enumerate}
\end{question}

\begin{theory}
For an \(M/M/1\) finite queue with maximum system size \(N\),
\[
p_n=p_0\rho^n,\qquad n=0,1,\ldots,N,
\]
and
\[
p_n=0,\qquad n>N.
\]

If
\[
\rho=1,
\]
then all finite-state probabilities are equal:
\[
p_0=p_1=\cdots=p_N=\frac{1}{N+1}.
\]

The blocking probability is
\[
P(\text{customer cannot enter})=p_N.
\]

The expected number in the system is
\[
L_s=\sum_{n=0}^{N}np_n.
\]
\end{theory}

\begin{answer}
Here,
\[
\lambda=4,\qquad \mu=4.
\]

Therefore,
\[
\rho=\frac{\lambda}{\mu}=1.
\]

Since \(N=4\), the possible states are
\[
0,1,2,3,4.
\]

\textbf{(a) Steady-state probabilities}

Because \(\rho=1\),
\[
p_0=p_1=p_2=p_3=p_4=\frac{1}{N+1}.
\]

\[
p_0=p_1=p_2=p_3=p_4=\frac{1}{5}.
\]

\[
\boxed{p_0=p_1=p_2=p_3=p_4=0.2.}
\]

\textbf{(b) Probability an arriving customer cannot enter}

A customer cannot enter when the system is full:
\[
P(\text{blocked})=p_4.
\]

\[
\boxed{P(\text{blocked})=p_4=\frac{1}{5}=0.2.}
\]

\textbf{(c) Expected number in the system}
\[
L_s=\sum_{n=0}^{4}np_n.
\]

\[
L_s=0\left(\frac{1}{5}\right)
+1\left(\frac{1}{5}\right)
+2\left(\frac{1}{5}\right)
+3\left(\frac{1}{5}\right)
+4\left(\frac{1}{5}\right).
\]

\[
L_s=\frac{0+1+2+3+4}{5}
=\frac{10}{5}=2.
\]

\[
\boxed{L_s=2.}
\]
\end{answer}

\section{State-Dependent Service Rate}

\begin{question}
A grocery store opens checkout counters depending on the number of customers in the store.
Customers arrive at rate
\[
\lambda_n=10\text{ per hour}.
\]
Each open counter serves at rate \(5\) customers per hour.
The number of open counters is:
\[
\begin{array}{c|c}
\text{Number of customers} & \text{Open counters}\\
\hline
1\text{ to }3 & 1\\
4\text{ to }6 & 2\\
7\text{ or more} & 3
\end{array}
\]
Write \(\mu_n\), the total service rate when there are \(n\) customers.
\end{question}

\begin{theory}
In a birth-death queue:
\[
\lambda_n=\text{arrival rate when in state }n,
\]
\[
\mu_n=\text{total departure rate when in state }n.
\]

If each server has service rate \(\mu\), and \(c_n\) servers are active in state \(n\), then
\[
\mu_n=c_n\mu.
\]
\end{theory}

\begin{answer}
Each open counter serves at rate
\[
5\text{ per hour}.
\]

If there are \(0\) customers, nobody can depart:
\[
\mu_0=0.
\]

If there are \(1,2,3\) customers, one counter is open:
\[
\mu_n=1(5)=5,\qquad n=1,2,3.
\]

If there are \(4,5,6\) customers, two counters are open:
\[
\mu_n=2(5)=10,\qquad n=4,5,6.
\]

If there are \(7\) or more customers, three counters are open:
\[
\mu_n=3(5)=15,\qquad n=7,8,9,\ldots
\]

Therefore,
\[
\boxed{
\mu_n=
\begin{cases}
0, & n=0,\\
5, & n=1,2,3,\\
10, & n=4,5,6,\\
15, & n\ge 7.
\end{cases}
}
\]
\end{answer}

\section{Kendall Notation}

\begin{question}
Explain the meaning of the queue notation
\[
(M/M/2):(FCFS/5/\infty).
\]
\end{question}

\begin{theory}
Kendall notation is
\[
(a/b/c):(d/e/f),
\]
where:
\[
a=\text{arrival distribution},
\]
\[
b=\text{service-time distribution},
\]
\[
c=\text{number of servers},
\]
\[
d=\text{queue discipline},
\]
\[
e=\text{maximum number allowed in the system},
\]
\[
f=\text{source size}.
\]

Here \(M\) means Markovian, meaning Poisson arrivals or exponential service times.
\end{theory}

\begin{answer}
For
\[
(M/M/2):(FCFS/5/\infty),
\]
we interpret each part:

\[
M=\text{Poisson arrival process}.
\]

\[
M=\text{exponential service time}.
\]

\[
2=\text{two parallel servers}.
\]

\[
FCFS=\text{first come, first served}.
\]

\[
5=\text{maximum of 5 customers in the system}.
\]

\[
\infty=\text{infinite calling population/source}.
\]

Thus this is a two-server queue with Poisson arrivals, exponential service times, first-come-first-served discipline, maximum system capacity \(5\), and unlimited possible customers in the source population.
\end{answer}

\section{Multiple-Server Queue Service Rate}

\begin{question}
An \(M/M/3\) queue has arrival rate \(\lambda=12\) customers per hour.
Each server serves at rate \(\mu=5\) customers per hour.
Write the total service rate \(\mu_n\) when there are \(n\) customers in the system.
\end{question}

\begin{theory}
For an \(M/M/c\) queue:
\[
\mu_n=
\begin{cases}
n\mu, & n<c,\\
c\mu, & n\ge c.
\end{cases}
\]

Reason:
\begin{itemize}
    \item if \(n<c\), only \(n\) servers are busy,
    \item if \(n\ge c\), all \(c\) servers are busy.
\end{itemize}
\end{theory}

\begin{answer}
Here,
\[
c=3,\qquad \mu=5.
\]

If \(n=0\), no one is in service:
\[
\mu_0=0.
\]

If \(n=1\), one server is busy:
\[
\mu_1=1(5)=5.
\]

If \(n=2\), two servers are busy:
\[
\mu_2=2(5)=10.
\]

If \(n\ge 3\), all three servers are busy:
\[
\mu_n=3(5)=15,\qquad n\ge 3.
\]

Therefore,
\[
\boxed{
\mu_n=
\begin{cases}
0, & n=0,\\
5, & n=1,\\
10, & n=2,\\
15, & n\ge 3.
\end{cases}
}
\]
\end{answer}

\section{Self-Service Model}

\begin{question}
In a self-service system, customers arrive at rate \(\lambda=8\) per hour.
Each customer completes service independently at rate \(\mu=2\) per hour.
There are infinitely many service positions, so no one waits.
Find the steady-state distribution \(p_n\).
\end{question}

\begin{theory}
For an infinite-server self-service model:
\[
\lambda_n=\lambda,
\]
and
\[
\mu_n=n\mu.
\]

Then
\[
p_n=\frac{\rho^n}{n!}p_0,
\qquad
\rho=\frac{\lambda}{\mu}.
\]

Using
\[
\sum_{n=0}^{\infty}p_n=1,
\]
we get
\[
p_0=e^{-\rho}.
\]

Therefore,
\[
p_n=e^{-\rho}\frac{\rho^n}{n!}.
\]
\end{theory}

\begin{answer}
Here,
\[
\lambda=8,\qquad \mu=2.
\]

So
\[
\rho=\frac{\lambda}{\mu}=\frac{8}{2}=4.
\]

For an infinite-server model,
\[
p_n=e^{-\rho}\frac{\rho^n}{n!}.
\]

Substitute \(\rho=4\):
\[
\boxed{p_n=e^{-4}\frac{4^n}{n!},\qquad n=0,1,2,\ldots}
\]

This is a Poisson distribution with mean \(4\).
\end{answer}

\section*{Essay-Type Revision Questions}

\begin{question}
Why is the condition \(\rho<1\) important in an \(M/M/1\) infinite queue?
\end{question}

\begin{answer}
For an \(M/M/1\) infinite queue,
\[
\rho=\frac{\lambda}{\mu}.
\]

Here \(\lambda\) is the arrival rate and \(\mu\) is the service rate.

If
\[
\rho<1,
\]
then
\[
\lambda<\mu.
\]

This means the server can serve customers faster than they arrive on average. The queue can settle into a steady state.

If
\[
\rho\ge 1,
\]
then arrivals are at least as fast as service. The queue tends to grow without bound, and a steady-state probability distribution does not exist for the infinite queue.
\end{answer}

\begin{question}
Explain the difference between \(L_s\) and \(L_q\).
\end{question}

\begin{answer}
\[
L_s=\text{expected number of customers in the whole system}.
\]

The system includes:
\[
\text{queue}+\text{service facility}.
\]

\[
L_q=\text{expected number of customers only waiting in the queue}.
\]

Therefore,
\[
L_s-L_q=\text{expected number of customers being served}.
\]

For a single-server queue, this value is also the probability that the server is busy.
\end{answer}

\end{document}
